%----------------------------------------------------------------------------------------
% DIAGONALISEREN. Alle opgaven zijn eigen werk.
%
% Opbouw: Probleem, Start, leerstof, en op een nieuwe pagina de Opdrachten, Koppel en Test jezelf.
% Labels van dit hoofdstuk beginnen met h3-.
%----------------------------------------------------------------------------------------

\mychapter{3}{Diagonaliseren}

\begin{probleem}
\naast{In 1202 stelde Leonardo van Pisa, bekend als Fibonacci, een raadsel over konijnen. Een pas geboren paar konijnen is na één maand volwassen, en een volwassen paar krijgt elke maand een nieuw paar. Er sterven geen konijnen. Start met één jong paar: het aantal paren is dan maand na maand
\[ 1,\ 1,\ 2,\ 3,\ 5,\ 8,\ 13,\ \ldots \]
want elke maand komen er zoveel jonge paren bij als er de maand ervoor paren waren. Hoeveel paren zijn er na 60 maanden, zonder alle tussenliggende termen te berekenen? En waarom gaat het quotiënt van twee opeenvolgende termen naar $1{,}618\ldots$?}{\begin{tikzpicture}[x=.5cm,y=.21cm,font=\scriptsize]
  \draw[black!14,very thin,ystep=5] (0,0) grid (8,22);
  \draw[->,grijs] (0,0) -- (8.6,0) node[right,text=black] {$n$};
  \draw[->,grijs] (0,0) -- (0,23.5) node[above,text=black] {paren};
  \foreach \t in {0,2,...,8} \node[below] at (\t,0) {\t};
  \foreach \t in {5,10,15,20} \node[left] at (0,\t) {\t};
  \foreach \k/\u in {0/1,1/1,2/2,3/3,4/5,5/8,6/13,7/21}{\fill[accent!55] ({\k-.3},0) rectangle ({\k+.3},\u); \draw[accent] ({\k-.3},0) rectangle ({\k+.3},\u);}
\end{tikzpicture}}
\end{probleem}

\lkr{Leerplan Complementair (III-WiCo-d): LPD 30, met de wenk over diagonaliseren: heeft een vierkante matrix een basis van eigenvectoren, dan kun je ze diagonaliseren; illustreer dat niet elke matrix zo'n basis heeft, en bereken machten van diagonaliseerbare matrices. Voorkennis: hoofdstuk 1 en 2, uit Matrices de inverse en het product, en uit Vectorruimten (of hoofdstuk 1) lineair onafhankelijk en basis: $n$ vectoren van $\R^n$ vormen een basis als en slechts als de matrix met die vectoren als kolommen een determinant verschillend van 0 heeft. Extra's: algebraïsche en meetkundige multipliciteit (de ongelijkheid zonder bewijs), diagonaliseren over $\C$, en de formule van Binet.}

In hoofdstuk 2 schreef je een begintoestand als combinatie van eigenvectoren, en kon je dan $A^nX_0$ berekenen zonder machten van $A$. In dit hoofdstuk doe je dat in één keer voor alle begintoestanden: je schrijft de matrix zelf met haar eigenvectoren en eigenwaarden. Dat heet \textit{diagonaliseren}. Het geeft een formule voor $A^n$, en het verklaart waarom de kolommen van een hoge macht van een overgangsmatrix gelijk worden. Maar het lukt niet bij elke matrix.

\begin{oefeningen}[Start]
\oef{h3-start-1}

\oef{h3-verken}
\end{oefeningen}

\lkr{De Verken toont het hoofdidee: met de eigenvectoren als kolommen van $P$ wordt $AP = PD$. Vraag c laat zien dat $P$ inverteerbaar is als de eigenvectoren een basis vormen; dat is de sleutel tot sectie 1.}

%========================================================================================
\section{Diagonaliseren}

\subsection{Een diagonaalmatrix}

Bij welke matrices is een macht eenvoudig? Bij een diagonaalmatrix: daar staan buiten de hoofddiagonaal alleen nullen, en we noteren ze $D = \diag(d_1, \ldots, d_n)$.

\begin{theorie}
\begin{stelling}{Machten van een diagonaalmatrix}
Voor elk natuurlijk getal $n$ is \qquad $\formule{\diag(d_1, \ldots, d_k)^n = \diag(d_1^n, \ldots, d_k^n)}$
\end{stelling}

\begin{bewijs}
Het product van twee diagonaalmatrices is de diagonaalmatrix met de producten op de hoofddiagonaal: in rij $i$ van de eerste staat alleen $d_i$ op plaats $i$, en in kolom $j$ van de tweede alleen $e_j$ op plaats $j$, dus het element $(i, j)$ van het product is $d_ie_i$ als $i = j$ en anders 0. Herhaal dat $n$ keer.
\end{bewijs}
\end{theorie}

De eenheidsvectoren zijn eigenvectoren van $D$: $D\vec e_i = d_i\vec e_i$. Een diagonaalmatrix rekt dus elke as uit met haar eigen factor. Een matrix met een basis van eigenvectoren doet hetzelfde, maar langs scheve assen. De figuur toont dat voor de matrix $A = \km{3 & -1 \\ 2 & 0}$ van hoofdstuk 1, met de eigenvectoren $\vec v_1 = (1, 1)$ bij 2 en $\vec v_2 = (1, 2)$ bij 1. In het scheve rooster van $\vec v_1$ en $\vec v_2$ is $A$ gewoon: verdubbel de eerste coördinaat, en laat de tweede staan.

\figuren{\begin{tikzpicture}[scale=.62,font=\scriptsize]
  \begin{scope}
    \scheefrooster{1}{1}{1}{2}{-1}{2}{-1}{1}
    \fill[accent!20] (0,0) -- (1,1) -- (2,3) -- (1,2) -- cycle;
    \vpijl{(0,0)}{(1,1)} \vpijl{(0,0)}{(1,2)}
    \node[below right] at (1,1) {$\vec v_1$};
    \node[left] at (1,2) {$\vec v_2$};
    \fill (0,0) circle (1.5pt);
  \end{scope}
  \draw[->,grijs,thick] (3.6,1) -- node[above] {$A$} (5.2,1);
  \begin{scope}[shift={(7.5,0)}]
    \scheefrooster{1}{1}{1}{2}{-1}{2}{-1}{1}
    \fill[fel!35] (0,0) -- (2,2) -- (3,4) -- (1,2) -- cycle;
    \bpijl{(0,0)}{(2,2)} \vpijl{(0,0)}{(1,2)}
    \node[below right] at (2,2) {$A\vec v_1 = 2\vec v_1$};
    \node[left] at (1,2) {$A\vec v_2 = \vec v_2$};
    \fill (0,0) circle (1.5pt);
  \end{scope}
\end{tikzpicture}}

\subsection{Een basis van eigenvectoren}

Hoe schrijf je dat in matrices? Zet de eigenvectoren als kolommen in een matrix $P$, en de eigenwaarden in dezelfde volgorde op de hoofddiagonaal van $D$. Dan is $AP$ de matrix met kolommen $A\vec v_1, A\vec v_2$, en $PD$ de matrix met kolommen $\lambda_1\vec v_1, \lambda_2\vec v_2$. Die zijn gelijk.

\begin{theorie}
Een vierkante matrix $A$ is \textbf{\textcolor{accent}{diagonaliseerbaar}} als er een inverteerbare matrix $P$ en een diagonaalmatrix $D$ bestaan met

\vspace{1mm}
\qquad $\formule{A = PDP^{-1}}$

\begin{stelling}{Diagonaliseerbaar}
Een $n \times n$-matrix $A$ is diagonaliseerbaar als en slechts als ze $n$ lineair onafhankelijke eigenvectoren heeft, dus een basis van eigenvectoren. De kolommen van $P$ zijn dan eigenvectoren, en de elementen van $D$ de eigenwaarden, in dezelfde volgorde.
\end{stelling}

\begin{bewijs}
Neem $P$ met kolommen $\vec v_1, \ldots, \vec v_n$ en $D = \diag(\lambda_1, \ldots, \lambda_n)$. Kolom $j$ van $AP$ is $A\vec v_j$, en kolom $j$ van $PD$ is $\lambda_j\vec v_j$, dus
\bewijsregel{AP = PD \iff A\vec v_j = \lambda_j\vec v_j \ \text{voor elke} \ j .}
Is $A = PDP^{-1}$, dan is $AP = PD$: elke kolom van $P$ is een eigenvector (niet $\nulv$, want $P$ is inverteerbaar), en de kolommen vormen een basis, want $\det P \neq 0$. Zijn omgekeerd $\vec v_1, \ldots, \vec v_n$ lineair onafhankelijke eigenvectoren, dan is $\det P \neq 0$, dus $P$ is inverteerbaar, en uit $AP = PD$ volgt $A = PDP^{-1}$.
\end{bewijs}
\end{theorie}

Waarom de naam? Uit $A = PDP^{-1}$ volgt $P^{-1}AP = D$: met de juiste basis wordt $A$ een diagonaalmatrix. Let op de volgorde: de eerste kolom van $P$ hoort bij het eerste element van $D$. Wissel je twee kolommen van $P$, dan moet je ook de eigenwaarden wisselen. En vermenigvuldig je een kolom van $P$ met een getal, dan blijft het een eigenvector: $P$ is dus niet uniek.

\voorbeeld{}{Diagonaliseer $B = \km{4 & -2 \\ 1 & 1}$ uit hoofdstuk 1.\label{vb:h3-diag}}{%
\uitleg{De eigenwaarden en eigenvectoren kennen we al. We zetten ze in $P$ en $D$, berekenen $P^{-1}$ en controleren.}
\begin{stappenbreed}
\lambda_1 = 2 \ \text{bij} \ (1, 1), \quad \lambda_2 = 3 \ \text{bij} \ (2, 1) & Voorbeeld~\ref{vb:h1-vec}. \\[1ex]
P = \km{1 & 2 \\ 1 & 1}, \quad D = \km{2 & 0 \\ 0 & 3} & De eigenvectoren als kolommen, de eigenwaarden in dezelfde volgorde. \\[3ex]
P^{-1} = \tfrac{1}{-1}\km{1 & -2 \\ -1 & 1} = \km{-1 & 2 \\ 1 & -1} & $\det P = 1 - 2 = -1 \neq 0$. \\[3ex]
PDP^{-1} = \km{2 & 6 \\ 2 & 3}\km{-1 & 2 \\ 1 & -1} = \km{4 & -2 \\ 1 & 1} & Controle: het klopt.
\end{stappenbreed}}

\begin{werkwijze}
Bereken de eigenwaarden en bij elke eigenwaarde een basis van de eigenruimte. Vind je samen $n$ lineair onafhankelijke eigenvectoren, zet ze dan als kolommen in $P$, en de eigenwaarden in dezelfde volgorde in $D$. Dan is $A = PDP^{-1}$. Controleer met $AP = PD$: dat vraagt geen inverse.
\end{werkwijze}

\nujij{Diagonaliseer $\km{1 & 4 \\ 1 & 1}$, met de eigenwaarden 3 bij $(2, 1)$ en $-1$ bij $(2, -1)$.}{$P = \km{2 & 2 \\ 1 & -1}$ en $D = \diag(3, -1)$. Controle: $AP = \km{6 & -2 \\ 3 & 1} = PD$.}

\begin{oefeningen}
\oef{h3-diag2}
\end{oefeningen}

\subsection{Verschillende eigenwaarden}

Moet je telkens nagaan of de eigenvectoren lineair onafhankelijk zijn? Niet als de eigenwaarden verschillen.

\begin{theorie}
\begin{stelling}{Verschillende eigenwaarden}
Eigenvectoren $\vec v_1$ en $\vec v_2$ bij verschillende eigenwaarden $\lambda_1 \neq \lambda_2$ zijn lineair onafhankelijk.
\end{stelling}

\begin{bewijs}
Stel $c_1\vec v_1 + c_2\vec v_2 = \nulv$. Vermenigvuldig met $A$, en ook met $\lambda_2$:
\bewijsregel{c_1\lambda_1\vec v_1 + c_2\lambda_2\vec v_2 = \nulv \quad \text{en} \quad c_1\lambda_2\vec v_1 + c_2\lambda_2\vec v_2 = \nulv .}
Trek ze af: $c_1(\lambda_1 - \lambda_2)\vec v_1 = \nulv$. Omdat $\lambda_1 \neq \lambda_2$ en $\vec v_1 \neq \nulv$, is $c_1 = 0$. Dan is $c_2\vec v_2 = \nulv$, dus ook $c_2 = 0$.
\end{bewijs}

\begin{stelling}{Genoeg verschillende eigenwaarden}
Heeft een $n \times n$-matrix $n$ verschillende reële eigenwaarden, dan is ze diagonaliseerbaar.
\end{stelling}

\begin{bewijs}
Neem bij elke eigenwaarde een eigenvector. Voor $n = 2$ zijn ze lineair onafhankelijk door de vorige stelling. Voor meer eigenvectoren werkt dezelfde redenering, telkens met één vector meer (opdracht~\ref{oef:h3-o21} doet $n = 3$). Met $n$ onafhankelijke eigenvectoren is $A$ diagonaliseerbaar.
\end{bewijs}
\end{theorie}

Het omgekeerde geldt niet: de eenheidsmatrix is al diagonaal, maar heeft maar één eigenwaarde. Met $n$ verschillende eigenwaarden hoef je dus niets na te gaan: diagonaliseren lukt altijd.

\voorbeeld{}{Diagonaliseer de matrix $C = \km{2 & 1 & -1 \\ 0 & 4 & 3 \\ 0 & -1 & 0}$ uit voorbeeld~\ref{vb:h1-drie}.\label{vb:h3-drie}}{%
\uitleg{De drie eigenwaarden 1, 2 en 3 verschillen: $C$ is diagonaliseerbaar. We controleren met $CP = PD$, kolom per kolom, zodat we $P^{-1}$ niet nodig hebben.}
\begin{stappenbreed}
P = \km{2 & 1 & -4 \\ -1 & 0 & -3 \\ 1 & 0 & 1}, \quad D = \km{1 & 0 & 0 \\ 0 & 2 & 0 \\ 0 & 0 & 3} & De eigenvectoren $(2, -1, 1)$, $(1, 0, 0)$ en $(-4, -3, 1)$. \\[5ex]
CP = \km{2 & 2 & -12 \\ -1 & 0 & -9 \\ 1 & 0 & 3} = PD & Elke kolom van $P$ maal haar eigenwaarde.
\end{stappenbreed}
Dus $C = PDP^{-1}$. De inverse $P^{-1}$ bereken je alleen als je ze nodig hebt, bijvoorbeeld voor $C^n$.}

\begin{oefeningen}
\oef{h3-diag3}
\end{oefeningen}

%========================================================================================
\section{Niet elke matrix is diagonaliseerbaar}

\subsection{Te weinig eigenvectoren}

Wat als een eigenwaarde meer dan één keer voorkomt? Neem de afschuiving $S = \km{1 & 1 \\ 0 & 1}$: ze schuift elk punt horizontaal op, evenveel als zijn hoogte. Haar enige eigenwaarde is 1, twee keer, en $S - I = \km{0 & 1 \\ 0 & 0}$ geeft $y = 0$: de eigenvectoren liggen allemaal op de $x$-as. Er is geen basis van eigenvectoren, dus $S$ is niet diagonaliseerbaar. Je ziet het ook zonder de stelling: was $S = PDP^{-1}$, dan stond in $D$ alleen de eigenwaarde 1, dus $D = I$ en $S = PIP^{-1} = I$. Maar $S \neq I$.

\figuren{\begin{tikzpicture}[scale=.75,font=\scriptsize]
  \rooster{-1}{1}{3}{0}{1}{2}
  \draw[grijs,thick] (0,0) rectangle (1,1);
  \fill[accent,opacity=.18] (0,0) -- (1,0) -- (2,1) -- (1,1) -- cycle;
  \draw[accent,very thick] (0,0) -- (1,0) -- (2,1) -- (1,1) -- cycle;
  \draw[accent!60,thick] (-1,0) -- (3.2,0) node[above left,text=accent] {$\Eig{1}$};
  \draw[->,fel!85!black,thick] (0,1.7) -- (1.7,1.7);
  \draw[->,fel!85!black,thick] (0,.85) -- (.85,.85);
\end{tikzpicture}}

Het verschil zit in twee getallen. Een eigenwaarde kan meerdere keren een nulwaarde zijn van de karakteristieke veelterm, en haar eigenruimte heeft een dimensie. Die twee hoeven niet gelijk te zijn.

\begin{theorie}
De \textbf{\textcolor{accent}{algebraïsche multipliciteit}} van een eigenwaarde $\lambda$ is het aantal keer dat $\lambda$ een nulwaarde is van $p_A$. De \textbf{\textcolor{accent}{meetkundige multipliciteit}} is $\formule{\dim \Eig{\lambda}}$, het aantal parameters van $(A - \lambda I)\vec v = \nulv$.

\begin{stelling}{Criterium}
Voor elke eigenwaarde geldt \quad $\formule{1 \leqslant \dim \Eig{\lambda} \leqslant \text{algebraïsche multipliciteit}}$

\vspace{1mm}
Een matrix met alleen reële eigenwaarden is diagonaliseerbaar als en slechts als bij elke eigenwaarde de twee multipliciteiten gelijk zijn.
\end{stelling}
\end{theorie}

Het bewijs van de ongelijkheid en van het criterium in het algemeen valt buiten deze cursus. Het idee: neem bij elke eigenwaarde een basis van de eigenruimte. Samen zijn die vectoren lineair onafhankelijk (zoals bij verschillende eigenwaarden), en hun aantal is de som van de meetkundige multipliciteiten. De som van de algebraïsche multipliciteiten is $n$, de graad van $p_A$. Je hebt dus $n$ eigenvectoren precies als nergens een eigenvector ontbreekt.

\voorbeeld{}{Is $N = \km{4 & 0 & 0 \\ 0 & 2 & 1 \\ -2 & 0 & 2}$ diagonaliseerbaar? En de matrix $D = \km{3 & 0 & 1 \\ 2 & 2 & 2 \\ 1 & 0 & 3}$ uit voorbeeld~\ref{vb:h1-dubbel}?\label{vb:h3-mult}}{%
\uitleg{Beide matrices hebben de eigenwaarden 2 (twee keer) en 4. Het verschil zit in de eigenruimte bij 2: we tellen de parameters. Bij een eigenwaarde die één keer voorkomt, zijn beide multipliciteiten altijd 1.}
\begin{stappenbreed}
N - 2I = \km{2 & 0 & 0 \\ 0 & 0 & 1 \\ -2 & 0 & 0} \colon\ x = 0, \ z = 0 & Eén parameter: $\dim \Eig{2} = 1 < 2$. \\[5ex]
N \ \text{niet diagonaliseerbaar} & Er ontbreekt een eigenvector. \\[1ex]
D - 2I = \km{1 & 0 & 1 \\ 2 & 0 & 2 \\ 1 & 0 & 1} \colon\ x + z = 0 & Twee parameters: $\dim \Eig{2} = 2$. \\[5ex]
P = \km{1 & 0 & 1 \\ 0 & 1 & 2 \\ -1 & 0 & 1}, \quad D' = \diag(2, 2, 4) & $D$ is wel diagonaliseerbaar, met $(1, 0, -1)$, $(0, 1, 0)$ en $(1, 2, 1)$.
\end{stappenbreed}}

\begin{werkwijze}
Bereken de eigenwaarden met hun algebraïsche multipliciteit. Zijn ze allemaal verschillend en reëel, dan is $A$ diagonaliseerbaar. Komt een eigenwaarde meer dan één keer voor, tel dan de parameters van $(A - \lambda I)\vec v = \nulv$. Is dat aantal kleiner dan de multipliciteit, dan is $A$ niet diagonaliseerbaar.
\end{werkwijze}

\nujij{Is $\km{5 & 0 \\ 3 & 5}$ diagonaliseerbaar?}{Eigenwaarde 5, twee keer. $A - 5I = \km{0 & 0 \\ 3 & 0}$ geeft $x = 0$: één parameter. Niet diagonaliseerbaar.}

\begin{oefeningen}
\oef{h3-niet}
\end{oefeningen}

\subsection{Complexe eigenwaarden}

De draaiing $R = \km{0 & -1 \\ 1 & 0}$ heeft geen reële eigenvectoren: over $\R$ is ze niet diagonaliseerbaar. Maar in hoofdstuk 1 vond je de complexe eigenwaarden $\pm i$, met eigenvectoren $(1, -i)$ en $(1, i)$. Die zijn lineair onafhankelijk, want hun eigenwaarden verschillen. Met complexe getallen gaat het dus wel:
\[ R = PDP^{-1} \quad \text{met} \quad P = \km{1 & 1 \\ -i & i} \quad \text{en} \quad D = \km{i & 0 \\ 0 & -i} . \]
Dat is meer dan een curiositeit. Zo is $R^4 = PD^4P^{-1} = PIP^{-1} = I$, omdat $i^4 = 1$: vier keer een kwartslag is een volle draai. In de opdrachten zie je dat ook de stellingen van dit hoofdstuk over $\C$ gelden. Wie alleen reële matrices $P$ en $D$ toelaat, moet zeggen: \textit{diagonaliseerbaar over $\R$}.

\begin{oefeningen}
\oef{h3-complex}
\end{oefeningen}

%========================================================================================
\section{Machten}

\subsection{Een formule voor \texorpdfstring{$A^n$}{A tot de n}}

Waarom is diagonaliseren nuttig? Omdat machten van $PDP^{-1}$ eenvoudig zijn: in het product $PDP^{-1} \cdot PDP^{-1}$ valt $P^{-1}P = I$ weg.

\begin{theorie}
\begin{stelling}{Machten van een diagonaliseerbare matrix}
Is $A = PDP^{-1}$, dan is voor elk natuurlijk getal $n$ \qquad $\formule{A^n = PD^nP^{-1}}$
\end{stelling}

\begin{bewijs}
Voor $n = 1$ staat het er. Stel dat $A^n = PD^nP^{-1}$. Dan is
\bewijsregel{A^{n+1} = A^n A = PD^nP^{-1} \cdot PDP^{-1} = PD^n(P^{-1}P)DP^{-1} = PD^{n+1}P^{-1} .}
Zo volgt de formule voor elke $n$, door inductie.
\end{bewijs}
\end{theorie}

Zo wordt een macht drie matrixproducten, wat $n$ ook is, en $D^n$ is gewoon een rij machten van getallen. Met letters krijg je een formule voor $A^n$.

\voorbeeld{}{Geef een formule voor $B^n$, met $B = \km{4 & -2 \\ 1 & 1}$ uit voorbeeld~\ref{vb:h3-diag}. Bereken $B^5$.\label{vb:h3-macht}}{%
\uitleg{We gebruiken $P$, $D$ en $P^{-1}$ uit voorbeeld~\ref{vb:h3-diag}, en vermenigvuldigen eerst $PD^n$.}
\begin{stappenbreed}
PD^n = \km{1 & 2 \\ 1 & 1}\km{2^n & 0 \\ 0 & 3^n} = \km{2^n & 2 \cdot 3^n \\ 2^n & 3^n} & Elke kolom van $P$ maal de macht van haar eigenwaarde. \\[3ex]
B^n = \km{2^n & 2 \cdot 3^n \\ 2^n & 3^n}\km{-1 & 2 \\ 1 & -1} = \km{2 \cdot 3^n - 2^n & 2^{n+1} - 2 \cdot 3^n \\ 3^n - 2^n & 2^{n+1} - 3^n} & Maal $P^{-1}$. Controle met $n = 1$: $B$. \\[3ex]
B^5 = \km{486 - 32 & 64 - 486 \\ 243 - 32 & 64 - 243} = \km{454 & -422 \\ 211 & -179} & Met $2^5 = 32$ en $3^5 = 243$.
\end{stappenbreed}}

\begin{werkwijze}
Diagonaliseer: $A = PDP^{-1}$. Bereken $P^{-1}$, dan $PD^n$ (elke kolom van $P$ maal de $n$-de macht van haar eigenwaarde), en ten slotte $PD^n \cdot P^{-1}$. Controleer met $n = 1$.
\end{werkwijze}

\begin{oefeningen}
\oef{h3-macht}
\end{oefeningen}

\subsection{Op lange termijn}

In de cursus Matrices zag je dat de kolommen van een hoge macht van een overgangsmatrix bijna gelijk worden, telkens aan het evenwicht. Met de formule voor $A^n$ zie je waarom.

\voorbeeld{}{Geef een formule voor $M^n$, met $M = \km{0{,}6 & 0{,}2 \\ 0{,}4 & 0{,}8}$ van de refter. Wat is de limiet van $M^n$?\label{vb:h3-markov}}{%
\uitleg{De eigenwaarden zijn 1 en $0{,}4$ (hoofdstuk 2), met eigenvectoren $(1, 2)$ en $(1, -1)$. We diagonaliseren en laten $n$ groot worden: $0{,}4^n$ gaat naar 0.}
\begin{stappenbreed}
P = \km{1 & 1 \\ 2 & -1}, \quad D = \km{1 & 0 \\ 0 & 0{,}4}, \quad P^{-1} = \tfrac13\km{1 & 1 \\ 2 & -1} & $\det P = -3$. \\[3ex]
M^n = \tfrac13\km{1 & 0{,}4^n \\ 2 & -0{,}4^n}\km{1 & 1 \\ 2 & -1} = \tfrac13\km{1 + 2 \cdot 0{,}4^n & 1 - 0{,}4^n \\ 2 - 2 \cdot 0{,}4^n & 2 + 0{,}4^n} & $PD^n$, maal $P^{-1}$. \\[3ex]
\lim_{n \to \infty} M^n = \tfrac13\km{1 & 1 \\ 2 & 2} & Elke kolom is het evenwicht $\bigl(\tfrac13, \tfrac23\bigr)$.
\end{stappenbreed}
De kolommen van $M^n$ worden dus gelijk omdat de tweede eigenwaarde in absolute waarde kleiner is dan 1: wat bij $0{,}4^n$ staat, verdwijnt. En omdat $M^nX_0$ een combinatie is van de kolommen met coëfficiënten van som 1, gaat elke begintoestand in aandelen naar hetzelfde evenwicht.}

\begin{oefeningen}
\oef{h3-markov}
\end{oefeningen}

\subsection{Terug naar het probleem}

Een rij waarin elke term afhangt van de twee vorige, zoals $u_{n+1} = u_n + u_{n-1}$, is een recursie. Met een vector maak je er een matrixmodel van, zoals in hoofdstuk 2. Zet twee opeenvolgende termen in een vector:
\[ X_n = \kol{u_{n+1} \\ u_n} \qquad \text{dan is} \qquad X_{n+1} = \kol{u_{n+1} + u_n \\ u_{n+1}} = \km{1 & 1 \\ 1 & 0}X_n . \]

\voorbeeld{}{Terug naar het probleem. In 1202 stelde Fibonacci een raadsel over konijnen. Een pas geboren paar konijnen is na één maand volwassen, en een volwassen paar krijgt elke maand een nieuw paar. Er sterven geen konijnen. Start met één jong paar. Het aantal paren $u_n$ na $n$ maanden is dan $1, 1, 2, 3, 5, 8, 13, \ldots$, met $u_0 = u_1 = 1$ en $u_{n+1} = u_n + u_{n-1}$: elke maand komen er zoveel jonge paren bij als er de maand ervoor paren waren. Geef een formule voor $u_n$, en bereken het aantal paren na 60 maanden zonder alle tussenliggende termen. Waarom gaat het quotiënt $\tfrac{u_{n+1}}{u_n}$ naar $1{,}618\ldots$?\label{vb:h3-terug}}{%
\uitleg{We diagonaliseren $F = \km{1 & 1 \\ 1 & 0}$. Omdat de eigenwaarden wortels bevatten, schrijven we ze met letters, en gebruiken we hun som en product. Daarna schrijven we $X_0 = (1, 1)$ als combinatie van de eigenvectoren, zoals in hoofdstuk 2.}
\begin{stappenbreed}
\lambda^2 - \lambda - 1 = 0 \colon\ \varphi = \tfrac{1 + \sqrt5}{2} \approx 1{,}618, \ \psi = \tfrac{1 - \sqrt5}{2} \approx -0{,}618 & Spoor 1, determinant $-1$: $\varphi + \psi = 1$, $\varphi\psi = -1$ en $\varphi - \psi = \sqrt5$. \\[2ex]
F - \varphi I = \km{1 - \varphi & 1 \\ 1 & -\varphi} \colon\ \vec v_1 = (\varphi, 1), \quad \vec v_2 = (\psi, 1) & De tweede rij: $x = \varphi y$. Zo ook bij $\psi$. \\[3ex]
(1, 1) = c_1(\varphi, 1) + c_2(\psi, 1) \colon\ c_1 = \tfrac{\varphi}{\sqrt5}, \ c_2 = -\tfrac{\psi}{\sqrt5} & $c_1 + c_2 = 1$ en $c_1\varphi + c_2\psi = 1$; gebruik $1 - \psi = \varphi$. \\[2ex]
u_n = c_1\varphi^n + c_2\psi^n = \dfrac{\varphi^{n+1} - \psi^{n+1}}{\sqrt5} & De tweede coördinaat van $X_n = c_1\varphi^n\vec v_1 + c_2\psi^n\vec v_2$. \\[3ex]
u_{60} = \dfrac{\varphi^{61} - \psi^{61}}{\sqrt5} = 2\,504\,730\,781\,961 & $\psi^{61}$ is verwaarloosbaar klein: $u_{60}$ is $\tfrac{\varphi^{61}}{\sqrt5}$, afgerond.
\end{stappenbreed}
Na 60 maanden zijn er dus ruim 2,5 biljoen paren. En omdat $\abs{\psi} < 1 < \varphi$, is $\varphi$ de dominante eigenwaarde: $u_n \approx \tfrac{\varphi^{n+1}}{\sqrt5}$, en dus gaat $\tfrac{u_{n+1}}{u_n}$ naar $\varphi \approx 1{,}618$, de \textit{gulden snede}.}

Een formule met $\sqrt5$ die voor elke $n$ een geheel getal geeft: dat is de \textit{formule van Binet} (1843), al kende De Moivre ze een eeuw eerder. Ze toont de kracht van diagonaliseren: de recursie wordt een som van twee meetkundige rijen, één per eigenwaarde.

\begin{oefeningen}
\oef{h3-recursie}
\end{oefeningen}

%========================================================================================
\opdrachtenkop

\begin{oefeningen}
\oef{h3-o1}

\oef{h3-o2}

\oef{h3-o3}

\oef{h3-o4}

\oef{h3-o5}

\oef{h3-o6}

\oef{h3-o7}

\oef{h3-o8}

\oef{h3-o9}

\oef{h3-o10}

\oef{h3-o11}

\oef{h3-o12}

\oef{h3-o13}

\oef{h3-o14}

\oef{h3-o15}

\oef{h3-o16}

\oef{h3-o17}

\oef{h3-o18}

\oef{h3-o19}

\oef{h3-o20}

\oef{h3-o21}
\end{oefeningen}

%----------------------------------------------------------------------------------------
% KOPPEL: verbanden met andere leerstof
%----------------------------------------------------------------------------------------
\needspace{20\baselineskip}
\opdrachtsectie{Koppel}

\lkr{Complexe getallen (de formule van De Moivre) en Goniometrische functies zijn gezien in pakket 4. Differentiaalvergelijkingen en Multivariate statistiek zijn modules van Complementair: geef die oefeningen alleen als de leerlingen ze zagen.}

\begin{oefeningen}
\oef{h3-koppel-1}

\oef{h3-koppel-2}

\oef{h3-koppel-3}
\end{oefeningen}

%----------------------------------------------------------------------------------------
\needspace{30\baselineskip} % de test en het schema Fout? Herhaal samen op een pagina
\opdrachtsectie{Test jezelf}

\begin{oefeningen}
\oef{h3-t1}

\oef{h3-t2}

\oef{h3-t3}

\oef{h3-t4}

\oef{h3-t5}

\oef{h3-t6}

\oef{h3-t7}
\end{oefeningen}

\herhaalschema{%
  \fv{oef:h3-t1}{vb:h3-diag}{oef:h3-diag2}%
  \fv{oef:h3-t2}{vb:h3-mult}{oef:h3-niet}%
  \fv{oef:h3-t3}{vb:h3-mult}{oef:h3-o9}%
  \fv{oef:h3-t4}{vb:h3-macht}{oef:h3-macht}%
  \fv{oef:h3-t5}{vb:h3-markov}{oef:h3-markov}%
  \fv{oef:h3-t6}{vb:h3-diag}{oef:h3-o2}%
  \fv{oef:h3-t7}{vb:h3-terug}{oef:h3-o11}}

\ifzelfstudie
\section*{Maak je eigen samenvatting}
Sluit dit boek en schrijf in je schrift: de machten van een diagonaalmatrix, wanneer een matrix diagonaliseerbaar is en hoe je $P$ en $D$ opstelt (met bewijs), waarom eigenvectoren bij verschillende eigenwaarden onafhankelijk zijn, een matrix die niet diagonaliseerbaar is en waarom, de twee multipliciteiten, de formule $A^n = PD^nP^{-1}$ met haar bewijs, en hoe je er een recursie zoals die van Fibonacci mee oplost. Vergelijk daarna met de kaders.
\fi
